\subsection{The class reciprocity pairing does not impose quintic parity}
\label{KT-QUINTIC-CLASS-PAIRING-SCOPE}

This is an applicability check, not a class-group vanishing theorem.
Let $p$ be odd, let $F=\mathbb Q(\mu_p)$, and write
$\Delta=\operatorname{Gal}(F/\mathbb Q)$ and $\omega$ for the mod-$p$
cyclotomic character.  The surviving component in the quintic question has
character $\alpha$ of order five and $\alpha(-1)=1$.  The previously proved
length bound allows this component to have order $p$.

\paragraph{Exact primary construction.}
The source is A.~Morgan and A.~Smith, \emph{Field change for the
Cassels--Tate pairing and applications to class groups}, Research in
Number Theory \textbf{10}, 61 (2024),
\texttt{doi:10.1007/s40993-024-00545-2}.
Its notation is
\[
 \operatorname{Cl}^{*}L=
 \operatorname{Hom}(\operatorname{Gal}(H_L/L),\mathbb Q/\mathbb Z).
\]
The perfect reciprocity pairing in (1.5), printed pp.~3--4, is between
$\operatorname{Cl}^{*}L[p^\infty]$ and
$\operatorname{Cl}L[p^\infty]$.  It is not an asserted alternating
self-pairing of the class group.

Definition~5.2, Notation~5.3, and Lemma~5.4, printed pp.~20--21, give
\[
 1\longrightarrow\mathcal O_L^\times/(\mathcal O_L^\times)^p
 \longrightarrow\operatorname{Sel}_pL
 \longrightarrow\operatorname{Cl}L[p]\longrightarrow1,
 \qquad
 \operatorname{Sel}_pL\simeq\operatorname{Sel}\mu_p.
 \tag{CP.1}
\]
Thus the Kummer Selmer input is not simply a second copy of
$\operatorname{Cl}L[p]$; it has the indicated unit kernel.
For the unramified exact sequence $E(a,b)$ in (5.3), Proposition~5.5
identifies its Cassels--Tate pairing with minus the reciprocity pairing
after this class-group projection.  The general pairing theorem recalled
as Theorem~2.1, printed p.~8, specifies kernels coming from the middle
Selmer group; it does not assert perfection before those kernels are
removed.

\paragraph{What the roots of unity supply.}
In \S5.2 the chosen identification of cyclic groups with roots of unity
gives a morphism $g_m:\frac1m\mathbb Z/\mathbb Z\to\mu_m$ in the
category of Selmer objects when $\mu_m\subset L$.  The inclusion of local
conditions used to define this morphism is not asserted to be an equality.
Proposition~5.6, printed p.~22, assumes $\mu_{ab}\subset L$ so that the
whole exact sequence, including its middle term, admits the requisite
map to the dual sequence.  It proves
\[
 \operatorname{CTP}_{E(a,b)}(\phi,g_a\psi)
 =\operatorname{CTP}_{E(b,a)}(\psi,g_b\phi).
 \tag{CP.2}
\]
For $a=b$ this is explicitly \emph{symmetry}.  In particular the case
$a=b=p$ requires $\mu_{p^2}\subset L$, which fails for
$L=\mathbb Q(\mu_p)$ and for its totally real quintic subfield.
The proof of Proposition~1.3 on the same page is an exact reformulation
of (CP.2) as a reciprocity identity; it does not change symmetry into
alternation.  A one-dimensional group $\mathbb Z/p\mathbb Z$ admits the
nondegenerate symmetric pairing $(x,y)\mapsto xy/p\bmod\mathbb Z$, so
symmetry alone cannot exclude the allowed order-$p$ case.

\paragraph{The actual character obstruction.}
The source defines Cartier duality by
$M^\vee=\operatorname{Hom}(M,(\overline{\mathbb Q})^\times)$
(\S2.1, printed p.~7).  Consequently
\[
 \bigl(\mathbb F_p(\alpha)\bigr)^\vee
   =\mathbb F_p(\omega\alpha^{-1}),\qquad
 \operatorname{Hom}_{G_{\mathbb Q}}\!
 \left(\mathbb F_p(\alpha),
       \mathbb F_p(\omega\alpha^{-1})\right)=0.
 \tag{CP.3}
\]
Indeed complex conjugation is $+1$ on the first character and $-1$ on
the second.  Equivalently, the required identity $\alpha^2=\omega$
is impossible by parity.  Even a proposed untwisted equivariant
self-duality with values in the constant module would require
$\alpha=\alpha^{-1}$, contrary to its order five.  The ordinary
perfect evaluation pairing pairs the character-$\alpha$ class piece
with its dual character piece in $\operatorname{Cl}^{*}$; it supplies
no identification of that piece with itself.

Passing to a field containing $\mu_{p^2}$ would change the class group
under discussion, and the proposition there still supplies the symmetric
identity (CP.2).  It does not manufacture the missing equivariant
self-duality of (CP.3) on the original component.  Hence the cited
construction supplies no perfect alternating self-pairing that would
force the possible single order-$p$ quintic component to vanish.
This conclusion is limited to the exact pairing hypotheses just checked;
it does not exclude other arithmetic input for the component.

\paragraph{Checked source scope.}
The complete requested material in \S\S5.1--5.2, including the proofs
of Lemma~5.4, Propositions~5.5--5.6 and the deduction of Proposition~1.3,
was read, together with the defining conventions and Theorem~2.1 in
\S2.  The general Cassels--Tate machinery from the authors' preceding
paper remains a cited input.  The published correction
\texttt{doi:10.1007/s40993-024-00552-3} adds an author affiliation and
contains no mathematical correction.
